% Generated draft: requires mathematical and source review.
\documentclass[11pt]{article}
\usepackage[margin=25mm]{geometry}
\usepackage{fontspec}
\setmainfont{Latin Modern Roman}
\usepackage{amsmath,amssymb,mathtools,unicode-math}
\setmathfont{Latin Modern Math}
\usepackage{enumitem,needspace,xurl,hyperref}
\setlength{\parindent}{0pt}
\setlength{\parskip}{5pt}
\setlength{\emergencystretch}{3em}
\newcommand{\N}{\mathbb{N}}
\newcommand{\Z}{\mathbb{Z}}
\newcommand{\Q}{\mathbb{Q}}
\newcommand{\R}{\mathbb{R}}
\newcommand{\C}{\mathbb{C}}
\newcommand{\F}{\mathbb{F}}
\newcommand{\A}{\mathbb{A}}
\newcommand{\PP}{\mathbb{P}}
\newcommand{\E}{\mathbb{E}}
\newcommand{\eps}{\varepsilon}
\newcommand{\dd}{\,\mathrm{d}}
\begin{document}
\section*{UnsolvedMath ID: 30001984}
\textbf{Source problem code:} catalogue entry
\subsection{Definitions and mathematical statement}
A subset of \(\N\) is computably enumerable (c.e.) if an algorithm enumerates its elements; this uses the usual natural-number coding convention. The supplied record does not define \(D\)-maximality, specify the action defining the orbits, identify which orbits are called ``remaining,'' or specify the reducibility intended by ``complete.'' Consequently, the following is a precise conditional formulation of its quantifiers, pending those definitions.

Let \(\mathscr O\) denote the intended family of orbits of \(D\)-maximal c.e. sets. The source describes \(\mathscr O\) as infinite. Let \(\mathscr R\subseteq\mathscr O\) be the subfamily referred to as the remaining orbits, and let \(\operatorname{Comp}(A)\) mean that \(A\) is complete in the intended, presently unspecified sense. The question is whether
\[
\exists F\subseteq\mathscr R\quad
\left[|F|<\infty\ \land\
\forall O\in\mathscr R\setminus F\;
\exists A\in O\;\operatorname{Comp}(A)\right].
\]
Equivalently, is the family
\[
\bigl\{O\in\mathscr R:\nexists A\in O\;\operatorname{Comp}(A)\bigr\}
\]
finite? The definitions of \(\mathscr O\), \(\mathscr R\), and \(\operatorname{Comp}\) must be recovered from the cited report before this becomes a fully specified mathematical problem. Infinitude of \(\mathscr O\) alone does not assert infinitude of \(\mathscr R\).
\subsection{Short English statement}
Must every remaining orbit of \(D\)-maximal computably enumerable sets, with at most finitely many exceptions, contain a complete set? The supplied record leaves the orbit action, completeness notion, definition of \(D\)-maximality, and scope of ``remaining'' unspecified.
\subsection{Sources}
\begin{itemize}
\item Local registry: \texttt{data/unsolvedmath/problems.json}, permanent UnsolvedMath ID \(30001984\), title ``Complete Sets in D-Maximal Orbits,'' source problem code \texttt{catalogue entry}. The supplied record labels the statement ``exact'' and its status ``open.'' No individual proposer is recorded; the proposed year is \(2012\).
\item Supplied report citation: \emph{Computability Theory} (2012), \emph{Oberwolfach Reports}, pp.~397--455, DOI \texttt{10.4171/owr/2012/08}. Source URL and supplied literature-source link, labelled ``original Oberwolfach report'': \url{https://doi.org/10.4171/owr/2012/08}. The background identifies Oberwolfach Reports as reports of mathematical workshops and their open-problem discussions.
\item Literature assessment dated exactly \texttt{2026-08-21}: the supplied curation says that the report leaves unknown whether the remaining \(D\)-maximal orbits contain complete sets and therefore classifies the extracted problem as open. This is supplied dated curation, not an independently verified assessment. The record also states that the mathematical content was checked against the official report and normalized without changing its source-level claim; that verification statement is supplied provenance and has not been independently checked here.
\end{itemize}
\end{document}
